\documentclass[11pt]{article}
\usepackage[margin=0.9in]{geometry}
\usepackage{amsmath,amssymb,amsthm,booktabs,mathrsfs}
\usepackage[T1]{fontenc}
\usepackage{lmodern,microtype}
\usepackage[colorlinks=true,linkcolor=blue,urlcolor=blue]{hyperref}
\newcommand{\norm}[1]{\left\lVert#1\right\rVert}
\newtheorem{theorem}{Theorem}[section]
\newtheorem{lemma}[theorem]{Lemma}
\newtheorem{proposition}[theorem]{Proposition}
\newtheorem{corollary}[theorem]{Corollary}
\theoremstyle{remark}
\newtheorem{remark}[theorem]{Remark}
\title{Toward a four-dimensional Yang--Mills construction\\
\large Notebook 37: the actual logarithm and finite local ultraviolet tuning}
\author{Research proof draft}
\date{17 October 2026}
\setlength{\emergencystretch}{5em}
\begin{document}
\maketitle
\begin{abstract}
We complete the leading-determinant comparison left open in Notebook 36,
under the explicitly stated classical structural and local quantum
inputs of this working series. The actual coefficient satisfies
\(\beta(\mathcal D_{M,K})=22\alpha_g\log K+O(1)\), uniformly in
preceding depth and in all applicable joint infinite structural systems.
Here \(\alpha_g=C_H/(48\pi^2)\) uses the original color normalization.

The proof has three new steps. First, source-localized sectorial
resolvents remove the soft infrared penalties through physical time
one with an integrable \(O(u)\) error. Second, nonlinear Ward identities
through field degree four compare the actual effective-action heat
coefficient with the Wilson coefficient with error \(O(s^{-1})\).
The left-relative chart-current terms are retained. Third, a direct
Wilson-to-continuum comparison gives an integrable \(O(s^{-1/2})\)
error and fixes the coefficient without importing a universality
assertion for an unmatched regulator.

The positive margin now supplies the previously missing hypothesis
of the finite local shooting theorem. It yields arbitrarily long
finite weak local histories with a fixed weak endpoint and bare
coupling of order the inverse square root of their length.
This is a theorem about the homogeneous local recursion.
Full Wilson measure matching across all field regions, sourced
continuum limits, and a nontrivial reconstructed Yang--Mills theory
remain unproved.
\end{abstract}

\section*{Scope and the proof dependency}

This is a research proof draft for mathematical review. Its new
arguments are given below; the previously constructed classical actions,
exact determinant identities, localized vertices and finite local
recursion are inputs from the earlier notebooks. Internal mathematical
audits and successful typesetting do not certify those upstream proofs.
The main result is not a completed four-dimensional existence theorem.
The deposited construction principally concerns \(SU(2)\).
General color-factor formulas below do not by themselves establish
all upstream construction inputs for every compact simple group.

Current coordinates use counting pairings. The primitive ratio
\(\ell>1\) is fixed, \(K=\ell^k\), and final physical spacing is
\(K^{-1}\). The preceding classical depth is \(M\).
The notation \([t^2]\) is a Taylor coefficient, whereas the Maxwell
functional is \(\beta=2[t^2]\) on the normalized affine probe.
For a Cartan element with \(\tau(H^2)=1\), put
\[
 C_H=\sum_{q\in\operatorname{spec}(\operatorname{ad}_H)}q^2,
 \qquad \alpha_g=C_H/(48\pi^2).
\]
Both signs of a charge are included in this adjoint trace; its
complexification introduces no extra multiplicity.

The infinite spatial estimates refer to joint structural systems
containing all required old-depth shifts and endpoint ratios. Their
constants are independent of the chosen system. Exact finite-period
determinant algebra and finite flat spectral bounds are distinguished
from continuous-momentum identities of these infinite kernels.
The finite local conclusion is recovered by the separate compactness
and moment-transfer theorem; no finite winding identity is inferred
from an infinite Fourier expansion.

The proof chain is
\[
 \begin{gathered}
 \text{Notebook 36: exact soft reduction and bounded heat ends}\\
 \Downarrow\\
 \text{soft penalty removal}
 \ \longrightarrow\ \text{actual base versus Wilson}\\
 \longrightarrow\ \text{Wilson versus continuum}
 \ \longrightarrow\ \text{positive actual logarithm}\\
 \longrightarrow\ \text{finite local transfer and shooting}.
 \end{gathered}
\]
The final section records the interfaces from this local theorem
to the original existence objective.



\section{Removing the soft penalties up to physical time one}
\label{spb-section}

Work in dimension four, in current counting coordinates, on the chosen
infinite structural system. Flat spatial bounds also pass to explicit periodizations by summing the exponentially decaying lifted kernels. This is not a linear-weight similarity on a finite torus. The affine coefficient itself is taken on the infinite structural system, not as a literal periodic magnetic flux; arbitrary actual finite winding kernels are not identified with those periodizations.
Set $\mu=K^{-1}$. The affine background below means only the first
two formal polynomial jets of
\[
 U_\alpha(t;x)=\exp(i t x_\nu H/K^2),\qquad \alpha\ne\nu.
\]
No nonzero unbounded field is inserted in a common analytic ball.

\subsection{Actual families and the required inputs}
\label{spb-inputs}

Use the actual effective classical Hessians and moving-reference maps:
\begin{align}
 B_s(A)&=D_A^TD_A,&
 B_v(A)&=S_M''(e^{i\mu A})+D_AD_A^T,\notag\\
 T_s(A)&=B_s(A)+\mu^2J_AJ_A^T,&
 \widetilde V(A)&=B_v(A)+\bar C_A^TF_A\bar C_A,\notag\\
 F_A&=aI-S_{MK}''(M_K(e^{i\mu A})).
 \label{spb-families}
\end{align}
All Hessians are left-relative, including their coordinate-current
terms. The auxiliary vector $\widetilde V$ has the already established
first two affine jets of the actual soft vector; equality off shell
is not asserted.

For either pair $(W,B)=(T_s,B_s)$ or
$(W,B)=(\widetilde V,B_v)$, put $X_E=K^2E$ and write
\begin{equation}
 X_W=X_B+C_A^TF_AC_A,
 \qquad
 (C_A,F_A)=
 \begin{cases}
 (J_A^T,I),&\text{scalar},\\
 (K\bar C_A,\ aI-S_{MK}''(M_K(e^{i\mu A}))),&\text{vector}.
 \end{cases}
 \label{spb-coarse-factor}
\end{equation}
The output fiber of $C_A$ has a fixed color/site or color/bond
dimension per final block. Choose the fixed $a\ge R_q+1$, so that
$F_0\ge I$ in the vector case.

The proof uses the following established inputs, with constants
uniform in $M,K$:

\begin{enumerate}
\item The actual flat $q_M$ has an exponential current-coordinate
kernel, analytic double zero, Maxwell lower bound and uniform second
moments. Hence, with the componentwise physical gradient
$\nabla=K d_0$ and $H_0=I+\nabla^T\nabla$,
\begin{equation}
 cH_0\le I+X_{B,0}\le CH_0.
 \label{spb-base-energy}
\end{equation}
\item The repaired covariant-Hessian theorem gives first-two physical
$C^1$ vertices, with complete smooth source collars and full rooted
operator bounds. The scaled gauge-generator products have the same
property.
\item The physical maps $K\bar C_A$, $J_A$ and their first-two
background derivatives have uniform whole-block operator bounds and
fixed physical range. The actual coarse Hessian has uniform rooted
source kernels. Consequently $C_A^TF_AC_A$ has uniformly bounded
zero-order physical vertices and full rooted source bounds.
\end{enumerate}
These inputs concern actual effective actions and corrected tangents.
No bare Wilson replacement, full uniform nonlinear physical domain,
or arbitrary high-frequency bounded-source estimate is used.

\begin{theorem}
\label{spb-main}
Under these inputs, if $P_0$ is the appropriate fine scalar or bond
projector onto final block zero, then
\begin{equation}
 \left|[t^2]\operatorname{Tr}
 P_0\bigl(e^{-uX_W(t)}-e^{-uX_B(t)}\bigr)\right|
       \le Cu,\qquad 0<u\le1,
 \label{spb-physical-bound}
\end{equation}
uniformly in old depth $M$ and endpoint $K$. Equivalently,
\begin{equation}
 \left|[t^2]\operatorname{Tr}
 P_0\bigl(e^{-sW(t)}-e^{-sB(t)}\bigr)\right|
       \le C\frac{s}{K^2},\qquad 0<s\le K^2.
 \label{spb-fine-bound}
\end{equation}
In particular the absolute $ds/s$ integral of the difference on this
entire band is bounded uniformly.
\end{theorem}

\subsection{Shifted flat resolvents with physical spatial weights}
\label{spb-resolvents}

At zero the penalty in \eqref{spb-coarse-factor} is nonnegative and
bounded uniformly. Thus \eqref{spb-base-energy} holds for $W$ too.
Fix $E=B$ or $W$ and set
\[
 \mathcal B_0=I+X_{E,0},\qquad R_z=(\mathcal B_0-z)^{-1}.
\]
Use a sectorial contour with rays $\arg z=\pm\pi/4$, $|z|\ge1/4$,
joined by the circular arc of radius $1/4$ to the left of the positive
spectrum. The spectral theorem and energy comparison give
\begin{align}
 \|R_z\|&\le C(1+|z|)^{-1},\notag\\
 \|\nabla R_z\|+\|R_z\nabla^T\|
      &\le C(1+|z|)^{-1/2},\notag\\
 \|\nabla R_z\nabla^T\|&\le C.
 \label{spb-energy-resolvent}
\end{align}
For example,
$\|(\mathcal B_0+|z|)^{1/2}R_z
(\mathcal B_0+|z|)^{1/2}\|\le C$ on this contour.

Only linear physical weights are needed for the nonlocal flat kernel.
Conjugate by $e^{\theta\cdot x/K}$, with sufficiently small real
$\theta$. The analytic double zero gives the fine-symbol estimate
\[
 \|q_M(p+i\mu\theta)-q_M(p)\|
 \le C\{\mu|\theta|\sqrt{\delta(p)}+\mu^2|\theta|^2\}.
\]
After multiplication by $K^2$, this is a small form perturbation
relative to $H_0$. The local gradient term has the same estimate.
The coarse correction also does: its physical legs have bounded norm
and fixed range, and its middle flat kernel has a uniform coarse
analytic strip. The estimate remains true relative to $H_0+|z|$.
A weighted resolvent Neumann series therefore proves
\eqref{spb-energy-resolvent} uniformly after conjugation.
Commuting a weight through $\nabla$ adds only $O(|\theta|)$ and
is covered by the scalar resolvent estimate.

For each pair of blocks choose the signs of $\theta$ along their
separation. The restricted versions of all three bounds then acquire
$e^{-m d(i,j)}$, for a common $m>0$. Internal block diameters change
only the constant. These are whole-block Hilbert operator estimates,
with no fine-entry trace or growing-fiber count. At fixed $K$ the
operators are bounded lattice operators; finite-support approximation
justifies the weighted kernels. All constants are uniform in $K$.

\subsection{At most one unbounded derivative per vertex}
\label{spb-vertices}

Let $X_j=[t^j]X_E(tA)$, $j=1,2$, for a complete physical $C^1$
direction of norm at most $a_*$. The repaired Hessian decomposition,
multiplied by $K^2$, gives
\begin{equation}
 X_j=\nabla^T\overline Z_j+\overline Z_j^T\nabla
                  +\overline E_j,
 \qquad
 \|\overline Z_j\|+\|\overline E_j\|\le C_j a_*^j.
 \label{spb-first-order}
\end{equation}
Indeed $\overline Z_j=KZ_j$ and $\overline E_j=K^2E_j$;
the improved $\overline Z_2=O(K^{-1})$ is not needed.
The scaled gauge-generator terms have this form, and the physical
coarse penalty contributes bounded zero-order terms. Mixed second
coefficients follow by polarization.

For directions in smooth physical collars $x,y$, each factor retains
an exponential connecting-tree bound including its operator endpoints
and every source. The significant fact is that every summand in
\eqref{spb-first-order} has at most one unbounded physical derivative.
An arbitrary second-order form bound would not give the following
integrable resolvent powers.

Only two jets are needed. If a full uniform nonlinear domain has not
been provided, use the polynomial family
$X(t)=X_0+tX_1+t^2X_2$, or its mixed quadratic version.
Equation \eqref{spb-first-order} gives a small common parameter disk
on which $I+X(t)$ is sectorial with the chosen contour. At each fixed
$K$ this is ordinary holomorphic functional calculus for bounded
operators. Its first-two semigroup jets equal the actual ones because
only $X_0,X_1,X_2$ enter their resolvent derivatives. Thus this device
adds no nonlinear-domain hypothesis.

\subsection{Rooted semigroup derivatives}
\label{spb-semigroup}

Differentiated resolvents satisfy
\begin{align}
 \delta_xR_z&=-R_zX_xR_z,\notag\\
 \delta_x\delta_yR_z
  &=R_zX_xR_zX_yR_z+R_zX_yR_zX_xR_z-R_zX_{xy}R_z.
 \label{spb-resolvent-derivatives}
\end{align}
One first-order insertion in \eqref{spb-first-order} costs
$(1+|z|)^{-3/2}$ by \eqref{spb-energy-resolvent}; its zero-order
part costs $(1+|z|)^{-2}$. Two first-order insertions cost at worst
$(1+|z|)^{-2}$. If both derivatives face the middle resolvent, that
piece costs $O(1)$ and both exterior resolvents cost
$(1+|z|)^{-1}$. In the other arrangements there are two half-power
energy legs and one full resolvent. The single $X_{xy}$ insertion
again has exponent $-3/2$.

Insert block projections between the already defined factors
$R_z\nabla^T$, $\nabla R_z$ and the bounded vertex factors.
No sharp block cutoff is commuted through a derivative. The weighted
resolvent and vertex bounds then give, after an exponent reserve,
\begin{align}
 \|\delta_x(R_z)_{ij}\|
  &\le C(1+|z|)^{-3/2}e^{-m\tau(i,j,x)},\notag\\
 \|\delta_x\delta_y(R_z)_{ij}\|
  &\le C\{(1+|z|)^{-3/2}+(1+|z|)^{-2}\}
                         e^{-m\tau(i,j,x,y)}.
 \label{spb-rooted-resolvent}
\end{align}
Every intermediate sum is over block operator norms; a smaller
exponent pays the finite convolution multiplicities.

Integrate these formulas on the Dunford contour for
$e^{-u\mathcal B_0}$. The rays have positive real part, the joining
arc has fixed length, and both powers in
\eqref{spb-rooted-resolvent} are integrable at infinity, uniformly
for $0<u\le1$. Since $e^{-uX}=e^u e^{-u(I+X)}$, one obtains
\begin{align}
 \|\delta_x(e^{-uX})_{ij}\|
     &\le Ce^{-m'\tau(i,j,x)},\notag\\
 \|\delta_x\delta_y(e^{-uX})_{ij}\|
     &\le Ce^{-m'\tau(i,j,x,y)}.
 \label{spb-rooted-semigroup}
\end{align}
For instance the ray integrals of $(1+\lambda)^{-3/2}$ and
$(1+\lambda)^{-2}$ are bounded by fixed constants even when the
factor $e^{-u\lambda/\sqrt2}$ is discarded.

The undifferentiated semigroup has a uniform exponential restricted
block bound too. The small linear-weight form perturbation retains
the positive gradient energy, leaving numerical range bounded below
by $-C$. Its real-time norm is at most $e^{Cu}$, and the weight gives
the block decay. This uses only the flat generator.

This resolvent argument avoids the nonintegrable middle-time factor
that an overly crude split of a second Duhamel insertion would produce.
It also retains the source roots which a global Cauchy bound alone
would lose.

\subsection{Affine reconstruction and the coarse trace}
\label{spb-trace}

Use a fixed smooth physical partition for the affine potential.
Its direction in collar $x$ has physical $C^1$ norm at most
$C(1+|x|)$. Summing its first and mixed second directions against
\eqref{spb-rooted-semigroup} gives, for $n=0,1,2$,
\begin{equation}
 \|([t^n]e^{-uX(t)})_{ij}\|
 \le C(1+|i|+|j|)^p e^{-m''d(i,j)},\qquad 0<u\le1.
 \label{spb-affine-semigroup}
\end{equation}
At $u=0$ use the identity and zero differentiated coefficients.
The sums are absolutely convergent polynomial-jet sums, not evaluations
at a nonzero unbounded field.

The physical $C_A,F_A$ in \eqref{spb-coarse-factor} have the same
polynomial block bounds; $C_A$ has fixed physical range and order-one
norm. In the vector case ordinary averaging sends the affine connection
to the coarse affine connection, so the actual coarse Hessian supplies
the required polynomial/exponential $F$ kernels.

Retain $P_0$ in the exact physical-time Duhamel trace identity:
\begin{equation}
 \operatorname{Tr}_fP_0(e^{-uX_W}-e^{-uX_B})
  =-\int_0^u\operatorname{Tr}_c
       [FCe^{-vX_B}P_0e^{-(u-v)X_W}C^T]\,dv.
 \label{spb-coarse-duhamel}
\end{equation}
The trace on the right sums all coarse diagonal blocks.
Its degree-two coefficient has finitely many ordered terms.
Equation \eqref{spb-affine-semigroup} and the $C,F$ estimates connect
both coarse endpoints to the retained root block. Their absolute
sum is bounded by a fixed multiple of
\[
 \sum_{Y,Z}(1+|Y|+|Z|)^p
            e^{-m'''(|Y|+|Z|+d(Y,Z))}<\infty.
\]
Each diagonal trace is in a fixed-dimensional coarse output fiber.
A product through the fine block-zero Hilbert space costs that coarse
rank, not its growing fine dimension. Thus the integrand coefficient
is bounded uniformly in $u,v,M,K$, including the endpoints $v=0,u$.
The integral has length $u$, proving \eqref{spb-physical-bound}.

For compact source truncations the cyclic trace step is ordinary
finite-rank cyclicity with $P_0$ retained. The displayed absolute
rooted sum gives the infinite limit and permits coefficient extraction.
This completes the proof of Theorem~\ref{spb-main}.

\subsection{Consequence and remaining comparison}
\label{spb-scope}

Equation \eqref{spb-fine-bound} implies
\begin{equation}
 \int_0^{K^2}\left|[t^2]\operatorname{Tr}
 P_0(e^{-sW(t)}-e^{-sB(t)})\right|\frac{ds}{s}\le C.
 \label{spb-integrated}
\end{equation}
Together with the coarse determinant reductions, this removes the
soft penalties from the fine-to-physical-unit cutoff heat coefficient
up to a uniform finite part. The vector conclusion is actual through
the established affine matching; the scalar family is actual directly.

The remaining middle-band comparison concerns the translation-invariant
actual base $S_M''+DD^T$ and the ordinary scalar Laplacian.
This section does not replace that effective Hessian by a fixed bare
reference, compute its logarithmic coefficient, or prove the separate
physical-time tail beyond one. A theorem for a fixed bare integrand
is not imported as a uniform statement in the present family.



\section{The complete mean-zero scalar coefficient}
\label{sgc-section}

The exact factorization of Notebook 36 reads
\[
 \det\widehat G=\det T_s\det M_s,\qquad
 T_s=D_A^TD_A+K^{-2}J_AJ_A^T.
\]
Here \(M_s=K^{-2}J_A^TT_s^{-1}J_A\) has one fixed color
fiber per final block. Its actual Maxwell coefficient is uniformly
bounded. The normalization is
\(\beta(F)=2[t^2]F\) on the specified one-color affine probe;
normalization at the flat field contributes zero.

\begin{theorem}[Mean-zero scalar logarithm]
\label{sgc-main}
Under the actual structural and affine-identification hypotheses of
Notebooks 35--36 and Theorem~\ref{spb-main},
\begin{equation}
 \left|\beta(\mathcal N_0\log\det\widehat G)
                 -2\alpha_g\log K\right|\le C,
 \qquad \alpha_g=\frac{C_H}{48\pi^2}.
 \label{sgc-coefficient}
\end{equation}
The bound is uniform in preceding depth and in every applicable
joint infinite structural system.
\end{theorem}

\begin{proof}
The scalar part of the infrared-tail theorem in Notebook 36 bounds
the integral over \(s>K^2\). Theorem~\ref{spb-main} removes the scalar
soft penalty on \(0<s\le K^2\) with a bounded integrated error.
The coarse factor \(M_s\) contributes a further bounded error.
The exact unpinned scalar heat calculation of Notebook 34 therefore
gives
\begin{equation}
 \beta(\mathcal N_0\log\det\widehat G)
       =\int_0^{K^2}h_g(s)\,\frac{ds}{s}+O(1),
 \qquad
 h_g(s)=\frac{C_H}{3}s^2e^{-8s}I_0(2s)^2I_1(2s)^2.
 \label{sgc-cutoff}
\end{equation}
The definition of \(h_g\) includes the factor \(-2\) multiplying
the scalar heat coefficient. The minus sign comes from the heat
representation of the logarithm, and the factor two from \(\beta\).
No additional real/complex color multiplicity is introduced.

The previously established elementary Bessel bounds imply
\[
 h_g(s)=O(s^4)\quad(0<s\le1),\qquad
 h_g(s)=\alpha_g+O(s^{-1})\quad(s\ge1).
 \]
Thus \(\int_0^1|h_g(s)|\,ds/s\) and
\(\int_1^\infty|h_g(s)-\alpha_g|\,ds/s\) are finite.
Splitting \eqref{sgc-cutoff} at one proves
\(\alpha_g\log(K^2)+O(1)\), including \(K=1\).
\end{proof}

This is the scalar determinant in the exact covariant mean-zero
decomposition. The earlier globally point-pinned scalar determinant
has a different logarithmic coefficient; its compensating vector
factor is different as well. Only each complete, exactly normalized
determinant combination is invariant under those changes of gauge
decomposition. This theorem alone does not compute the vector
coefficient or prove positive total drift.


% No preamble. Prefix abw-. Audited research, 17 October 2026.
\section{Uniform actual-base heat comparison from nonlinear Ward jets}
\label{abw-section}

This section concerns every applicable joint infinite structural system
with the established homogeneous classical inputs below. Constants are
uniform in the old classical depth \(M\). We do not interpolate a finite
momentum grid or identify finite winding kernels with periodizations.

Let \(S_M\) be the actual homogeneous classical minimum action in current
unit coordinates, with a fixed \(4\dim\mathfrak g\) input fiber.
Use its translation, hypercubic, reality and nonlinear gauge Ward
identities, and its uniformly exponentially rooted kernels through field
order four. The needed finite spatial moments, local exponential-chart
derivatives and their constants are uniform in \(M\).
Write \(S_W=S_1\) for the normalized Wilson action. The established flat
normalization and lower bound are
\begin{equation}
 q_M\ge q_W=d_{\rm curl}^{T}d_{\rm curl},\qquad q_MD=0,\qquad
 q_{M,\mathrm{mid}}(p)=|p|^2I-pp^T+O(|p|^4).
 \label{abw-input}
\end{equation}
These are the actual classical current-unit bounds. A bound for an
arbitrary quantum activity with constants depending on a newly chosen
auxiliary grouping would not replace them.

The background-Feynman vector base, in the actual left-relative link
chart, is
\[
 {\cal B}_M(a)=S_M''(\exp(ia))+D_aD_a^T.
\]
For \(\tau(H^2)=1\) and distinct directions \(\alpha,\nu\), define its
affine heat coefficient by its convergent fixed-time Taylor kernels:
\begin{equation}
 h_M(s)=[t^2]\operatorname{tr}_{E}
 \bigl(e^{-s{\cal B}_M(a(t))}\bigr)_{00},\qquad
 a_\alpha(t,x)=t x_\nu H.
 \label{abw-h}
\end{equation}
This is an infinite-lattice density, not a varying-flux family on a
finite torus.

\begin{theorem}[Uniform base comparison]
\label{abw-main}
Under these inputs there is \(C\), independent of \(M\), such that
\begin{equation}
 |h_M(s)-h_W(s)|\le C s^{-1}\quad(s\ge1),\qquad
 \int_1^{K^2}|h_M(s)-h_W(s)|\,\frac{ds}{s}\le C.
 \label{abw-bound}
\end{equation}
The scalar base \(D_a^TD_a\) is identical for \(M\) and \(W\).
Thus the same comparison applies to the vector-minus-two-scalar
combination. Explicit periodizations may also be used where their
identification has been established.
\end{theorem}

\subsection{Actual relative vertices, including chart curvature}
\label{abw-chart-section}

Set \(R_M=S_M-S_W\). Its Maxwell coefficient vanishes by
\eqref{abw-input}. The previously proved total nonlinear Ward extraction
therefore gives, with uniform analytic moment bounds,
\begin{equation}
 R_2=O(|p|^4),\qquad R_3=O((\max_i|p_i|)^3),\qquad
 R_4=O((\max_i|p_i|)^2).
 \label{abw-ward-orders}
\end{equation}
Analytic changes between midpoint and source-site bond phases preserve
these orders.

For \(\Xi_a(B)=\log(e^{iB}e^{ia})/i\), write
\[
 J(a)=D_B\Xi_a(0)=I+J_1(a)+J_2(a,a)+O(a^3),\qquad
 T(a)=D_B^2\Xi_a(0)=T_1(a)+O(a^2).
\]
In particular \(T(0)=0\). The exact relative Hessian is
\begin{equation}
 L_R(a)[B,C]=R_M''(a)[J(a)B,J(a)C]
                  +R_M'(a)[T(a)[B,C]].
 \label{abw-chain}
\end{equation}
At background degree one it is
\[
 R_3[a,B,C]+R_2[J_1(a)B,C]+R_2[B,J_1(a)C].
\]
Its degree-two coefficient is
\begin{align}
 &\tfrac12R_4[a,a,B,C]
 +R_3[a,J_1(a)B,C]+R_3[a,B,J_1(a)C]\notag\\
 &\quad+R_2[J_2(a,a)B,C]+R_2[B,J_2(a,a)C]
 +R_2[J_1(a)B,J_1(a)C]+R_2[a,T_1(a)[B,C]].
 \label{abw-second-chart}
\end{align}
Mixed derivatives are obtained by polarization. The last term is the
current times chart-curvature term and is retained.

All coefficients of \(J_1,J_2,T_1\) are local at one bond with fixed color
norms. They only merge momenta and introduce no inverse momentum.
Thus the flat, first-background and second-background differences of the
actual Hessians have respective momentum orders \(4,3,2\).
The gauge-fixing term \(D_aD_a^T\) is the same for both actions and
cancels from these differences.

\subsection{Fourier bounds and exact total heat variations}
\label{abw-diagram-section}

Put \(B_M(k)={\cal B}_M(0,k)\) and
\(\delta(k)=4\sum_\mu\sin^2(k_\mu/2)\). Then
\begin{equation}
 B_W(k)=\delta(k)I,\qquad
 c\delta(k)I\le B_M(k)\le C\delta(k)I.
 \label{abw-ellipticity}
\end{equation}
The upper bound follows from the uniform quadratic zero and bounded
second moments. On a fixed small punctured ball, with \(\rho=|k|\),
\begin{equation}
 \|\partial^r(B_M-B_W)\|\le C\rho^{4-r}\ (r\le2),\qquad
 \|\partial B_M\|\le C\rho,\qquad \|\partial^2B_M\|\le C.
 \label{abw-flat-jets}
\end{equation}

Let \(V_a^M(k,p)\) be the first vertex for a background
\(e^{ipx}e_a\), mapping quantum momentum \(k\) to \(k+p\).
Let \(W_{ab}^M(k;p,-p)\) be the mixed second vertex, preserving \(k\).
For \(\Delta=M-W\), \eqref{abw-ward-orders}--\eqref{abw-second-chart}
imply at \(p=0\), for \(r\le2\),
\begin{align}
 \|\partial_p^r\Delta V_a(k,p)\|
 +\|\partial_p^r\Delta V_b(k+p,-p)\|
       &\le C\rho^{3-r},\notag\\
 \|\partial_p^r\Delta W_{ab}(k;p,-p)\|
       &\le C\rho^{2-r}.
 \label{abw-vertex-difference}
\end{align}
The second first-vertex term includes differentiation of both arguments.
For either baseline it suffices to use
\begin{equation}
 \|\partial_p^rV\|\le C\rho^{1-r},\qquad
 \|\partial_p^rW\|\le C\rho^{-r}.
 \label{abw-baseline-vertices}
\end{equation}
Negative powers here are only loose bounds on analytic functions.
In particular \(\partial_p^2W\) is bounded. Away from the small ball all
required vertices and derivatives are uniformly bounded.

For compact source variations the second total heat variation is
well-defined by trace-class Duhamel insertions. There is no need to assign
a value to the extensive vacuum trace. Differentiating
\(D\operatorname{Tr}e^{-sB}[a]=-s\operatorname{Tr}(e^{-sB}B'[a])\)
gives
\begin{equation}
 D^2\operatorname{Tr}e^{-sB}[a,b]
 =-s\operatorname{Tr}(e^{-sB}B''[a,b])
 +s^2\int_0^1\operatorname{Tr}
  (e^{-s(1-v)B}B'[a]e^{-svB}B'[b])\,dv.
 \label{abw-total-duhamel}
\end{equation}
The integral already includes both orders; there is no further factor
two. Only the \emph{total} trace was cycled.

The external-momentum kernel is accordingly
\begin{align}
 {\cal H}_{ab}^M(s,p)
 &=-s\int_{\mathrm{BZ}}\operatorname{tr}
       (e^{-sB_M(k)}W_{ab}^M(k;p,-p))\,\frac{dk}{(2\pi)^4}
       \notag\\
 &\quad+s^2\int_0^1\int_{\mathrm{BZ}}\operatorname{tr}
 \left(e^{-s(1-v)B_M(k+p)}V_a^M(k,p)
       e^{-svB_M(k)}V_b^M(k+p,-p)\right)
                          \frac{dk}{(2\pi)^4}\,dv.
 \label{abw-diagram}
\end{align}
At fixed \(s\) the kernels have the moments permitting two external
derivatives under this integral.

\subsection{Endpoint-safe heat estimates and the gain}
\label{abw-estimate-section}

Duhamel differentiation and \eqref{abw-ellipticity} give, for
\(E_M(t,k)=e^{-tB_M(k)}\) and \(r=0,1,2\),
\begin{equation}
 \|\partial_k^rE_M(t,k)\|
 \le C_r\rho^{-r}(1+t\rho^2)^r e^{-ct\rho^2}.
 \label{abw-heat}
\end{equation}
The sharper first and second derivative factors are respectively
\(Ct\rho\) and \(C(t+t^2\rho^2)\). No spectral frame is differentiated.
The exact difference identity
\[
 E_M(t)-E_W(t)
 =-\int_0^t E_M(t-u)(B_M-B_W)E_W(u)\,du
\]
and \eqref{abw-flat-jets} give
\begin{equation}
 \|\partial_k^r(E_M-E_W)(t,k)\|
 \le C_r t\rho^{4-r}(1+t\rho^2)^r e^{-ct\rho^2}.
 \label{abw-heat-difference}
\end{equation}
These bounds hold for every \(t\ge0\); in particular they contain no
singularity at either endpoint \(v=0,1\) of \eqref{abw-diagram}.

Differentiate \eqref{abw-diagram} twice externally and subtract.
For the tadpole, the two differences are
\((E_M-E_W)\partial_p^2W_W\) and
\(E_M\partial_p^2(W_M-W_W)\).
Before the outer factor \(s\), they are bounded by
\(Cs\rho^4e^{-cs\rho^2}\) and \(Ce^{-cs\rho^2}\).
Their integrals are \(O(s^{-2})+O(s^{-1})\) after that factor.

Telescope the bubble product into its four factor differences.
Replacing a vertex gives combined vertex degree \(3+1=4\).
After distributing two derivatives between vertices and heat factors,
the full product is bounded by
\[
 C\rho^2 P(s\rho^2)e^{-cs\rho^2},
\]
where \(P\) is a fixed polynomial. Replacing a heat factor instead
gives, by \eqref{abw-heat-difference}, a time at most \(s\) and degree
\(4+1+1\); the differentiated bound is
\[
 Cs\rho^4 P(s\rho^2)e^{-cs\rho^2}.
\]
Both bounds are uniform in \(v\), since the two heat times sum to \(s\).
The four-dimensional integrals of these expressions are
\(O(s^{-3})\) and \(O(s^{-3})\), respectively.
The prefactor \(s^2\) therefore gives \(O(s^{-1})\).
Away from zero, a fixed spectral gap gives a polynomial in \(s\)
times \(e^{-cs}\). We have proved
\begin{equation}
 |\partial_p^2({\cal H}_{ab}^M-{\cal H}_{ab}^W)(s,0)|
 \le C/s,\qquad s\ge1.
 \label{abw-two-point-bound}
\end{equation}

\subsection{Density normalization and scope}
\label{abw-density-section}

The diagonal color heat density is gauge invariant, and its linear
color jet is zero for a semisimple group. The established density
Ward/reanchoring identity at fixed \(s\) consequently gives
\begin{equation}
 h_M(s)=\tfrac14\,
 \partial_{p_\nu}^2{\cal H}_{\alpha H,\alpha H}^M(s,0).
 \label{abw-density-factor}
\end{equation}
Indeed the Maxwell extraction is both twice the affine quadratic
density and one half of this total Hessian derivative. On the diagonal
component the source-site and midpoint phases cancel.

For the actual endpoint-\(K\) source \(a_\alpha=t x_\nu H/K^2\),
translation changes the background only by a constant pure gauge.
Every diagonal coefficient in a final block is therefore the same.
The \(K^4\) sites cancel the \(K^{-4}\) from the two source factors.
Its block coefficient is exactly \eqref{abw-h}; no partial-trace
cyclicity or finite-torus flux interpolation enters.
Equations \eqref{abw-two-point-bound} and \eqref{abw-density-factor}
prove Theorem~\ref{abw-main}.

This result supplies the uniform old-depth comparison for the base
operator. The soft-penalty, infrared, short-time and coarse-determinant
comparisons remain separate inputs to an actual determinant assembly.
The proof does not by itself construct a finite-regulator running
trajectory, a large-field RG history, a continuum measure or a mass gap.



% No preamble. Prefix wcr-. Fixed reference, 17 October 2026.
\section{Direct Wilson-to-continuum heat comparison}
\label{wcr-section}

We now evaluate the fixed reference by its actual vertices and diagonal
heat kernels. No general one-loop universality theorem is imported.
Let \(h_{W,v}(s)\) and \(h_{W,g}(s)\) be the affine Cartan coefficients
of the Wilson background-Feynman vector base and the full scalar
covariant Laplacian on \(\mathbb Z^4\). Normalize
\[
 \tau(H^2)=1,\qquad C_H=\operatorname{Tr}_{\rm ad}
       (\operatorname{ad}_H^2),\qquad \alpha_g=\frac{C_H}{48\pi^2}.
\]

\begin{theorem}[Fixed-reference heat coefficient]
\label{wcr-main}
For \(s\ge1\),
\begin{equation}
 |h_{W,v}(s)-10\alpha_g|
 +|h_{W,g}(s)+\alpha_g/2|\le Cs^{-1/2}.
 \label{wcr-fixed-bound}
\end{equation}
Together with the actual-base comparison, this gives on every
applicable joint infinite structural system
\begin{equation}
 h_{M,v}(s)-2h_g(s)=11\alpha_g+O(s^{-1/2}),\qquad
 \int_1^{K^2}\bigl(h_{M,v}(s)-2h_g(s)\bigr)\frac{ds}{s}
       =22\alpha_g\log K+O(1),
 \label{wcr-middle}
\end{equation}
uniformly in the old classical depth \(M\).
\end{theorem}

\subsection{The exact Wilson density fixes the leading vertices}
\label{wcr-density-section}

The normalized unit-lattice action is
\[
 S_W(U)=\sum_{x,\mu<\nu}
 \left[1-\tfrac12\tau(P_{x,\mu,\nu}+P_{x,\mu,\nu}^{-1})\right].
\]
Each unoriented plaquette is counted once. Its left-relative Hessian and
scalar covariant gradient define
\[
 {\cal B}_W(a)=S_W''(e^{ia})+D_aD_a^T,\qquad
 L_W(a)=D_a^TD_a,\qquad
 (D_a\phi)_\mu(x)=\operatorname{Ad}_{e^{ia_\mu(x)}}\phi(x+e_\mu)-\phi(x).
\]
Transposes use the fixed bilinear counting pairing.

Introduce a physical mesh \(\epsilon\), with link angle \(\epsilon A\).
For the physical jet
\[
 z=(a,b,c,e)
 =\bigl(A_\mu,A_\nu,\nabla_{\epsilon,\mu}A_\nu,
                         \nabla_{\epsilon,\nu}A_\mu\bigr)
\]
at one plaquette, its holonomy is exactly
\begin{equation}
 P_\epsilon(z)=e^{i\epsilon a}e^{i\epsilon b+i\epsilon^2c}
              e^{-i\epsilon a-i\epsilon^2e}e^{-i\epsilon b}.
 \label{wcr-exact-plaquette}
\end{equation}
Since \(P_0=I\) and \(\partial_\epsilon P_0=0\),
\[
 F_\epsilon=\frac{P_\epsilon-I}{\epsilon^2},\qquad
 {\cal L}_\epsilon=-\tfrac12\tau(F_\epsilon^2P_\epsilon^{-1})
\]
extend jointly holomorphically to \(\epsilon=0\).
The identity \(P+P^{-1}-2I=(P-I)^2P^{-1}\) gives the exact action
\begin{equation}
 S_\epsilon(A)=\epsilon^4\sum_{x,\mu<\nu}
                 {\cal L}_\epsilon(J_\epsilon A).
 \label{wcr-exact-density}
\end{equation}
Direct differentiation of \eqref{wcr-exact-plaquette} gives
\begin{equation}
 F_0=i(c-e)-[a,b],\qquad
 {\cal L}_0=\tfrac12\tau\bigl(c-e+i[a,b]\bigr)^2.
 \label{wcr-density-limit}
\end{equation}
On each fixed bounded complex jet polydisc,
\({\cal L}_\epsilon-{\cal L}_0\), together with every fixed field
derivative, is \(O(\epsilon)\). Joint holomorphy and Cauchy's formula
supply this estimate before any momentum limit is taken.

On plane waves with bounded complex physical momenta,
\((e^{i\epsilon p_\mu}-1)/\epsilon=ip_\mu+O(\epsilon)\)
jointly analytically. Equations \eqref{wcr-exact-density} and
\eqref{wcr-density-limit} therefore identify the continuum action as
\begin{equation}
 S_c(A)=\tfrac12\int\sum_{\mu<\nu}\tau(F_{\mu\nu}(A)^2)\,dx,\qquad
 F_{\mu\nu}=\partial_\mu A_\nu-\partial_\nu A_\mu+i[A_\mu,A_\nu].
 \label{wcr-continuum-action}
\end{equation}
In particular, for the ordinary action vertices,
\begin{equation}
 \Gamma_n^W(\epsilon p_1,\ldots,\epsilon p_n)
 =\epsilon^{4-n}\Gamma_n^c(p_1,\ldots,p_n)
                  +O(\epsilon^{5-n}),\qquad n=2,3,4.
 \label{wcr-engineering}
\end{equation}
This follows by differentiating the exact density in plane-wave
amplitudes; it holds jointly with the fixed momentum derivatives used
below. In midpoint coordinates the bare quadratic correction is fourth
order. Source-site phases may give third-order bare cross entries.
The full flat Feynman symbol is exactly \(\delta(k)I\) and has the
fourth-order correction needed below, in either convention.

The physical left-relative chart is
\[
 \Xi_{\epsilon,A}(B)
   =\frac{\log(e^{i\epsilon B}e^{i\epsilon A})}{i\epsilon}
   =A+B+O(\epsilon).
\]
Its first \(B\)-derivative is \(I+O(\epsilon)\), its second is
\(O(\epsilon)\). The exact relative Hessian includes both terms
\[
 S''(A)[J_AB,J_AC]+S'(A)[T_A(B,C)].
\]
Thus it has the same continuum limit; the current term is retained.
At first background order its local corrections multiply
\(\Gamma_2\), hence have momentum order at least two.
At second background order the corrections have order at least one
from \(\Gamma_3\), or two from \(\Gamma_2\) and the current term.
They are within the remainders in \eqref{wcr-engineering}.

Likewise,
\[
 \epsilon^{-1}D_{\epsilon A}
 =\epsilon^{-1}
   \bigl(\operatorname{Ad}_{e^{i\epsilon A_\mu}}
                          \mathrm{shift}_{\epsilon e_\mu}-I\bigr)
 \longrightarrow D_{c,A,\mu}=\partial_\mu+i\operatorname{ad}_{A_\mu}
\]
jointly on these jets. Integration by parts in
\eqref{wcr-continuum-action} shows that the limiting vector and scalar
operators are
\begin{equation}
 ({\cal B}_c(A))_{\mu\nu}
 =-D_{c,A}^2\delta_{\mu\nu}-2i\operatorname{ad}_{F_{\mu\nu}},
 \qquad L_c(A)=-D_{c,A}^2.
 \label{wcr-continuum-operators}
\end{equation}
The opposite consistent adjoint sign convention gives identical heat
traces. The coefficient of the flat Laplacian is one.

\subsection{An integrable comparison of the two-point diagrams}
\label{wcr-comparison-section}

Use the same source-site Fourier convention for lattice and continuum.
For either scalar or vector their flat symbols are
\[
 B_W(k)=\delta(k)I,\qquad B_c(k)=|k|^2I,\qquad
 \delta(k)=4\sum_\mu\sin^2(k_\mu/2).
\]
On a fixed small ball with \(\rho=|k|\), their difference and its
derivatives through order two are \(O(\rho^{4-r})\).
Define \(V_a(k,p)\) and \(W_{ab}(k;p,-p)\) as the first and second
background vertices, as in the actual-base comparison.
The joint analytic expansion above gives
\begin{align}
 \|\partial_p^r(V_W-V_c)(k,p)|_{p=0}\|
 +\|\partial_p^r(V_W-V_c)(k+p,-p)|_{p=0}\|
     &\le C\rho^{2-r},\notag\\
 \|\partial_p^r(W_W-W_c)(k;p,-p)|_{p=0}\|
     &\le C\rho^{1-r},\qquad r\le2.
 \label{wcr-vertex-errors}
\end{align}
Here \(V_c\) is homogeneous of degree one and \(W_c\) is constant.
The \(r=2\) bound \(C/\rho\) on the second line is deliberately loose:
the corresponding lattice derivative is analytic and bounded.
For either baseline we can use
\(\|\partial_p^rV\|\le C\rho^{1-r}\) and
\(\|\partial_p^rW\|\le C\rho^{-r}\).

The exact total heat Hessian kernel is
\begin{align}
 {\cal H}_{ab}(s,p)
 &=-s\int\operatorname{tr}
       (e^{-sB(k)}W_{ab}(k;p,-p))\,\frac{dk}{(2\pi)^4}
       \notag\\
 &\quad+s^2\int_0^1\int\operatorname{tr}
 \left(e^{-s(1-v)B(k+p)}V_a(k,p)
       e^{-svB(k)}V_b(k+p,-p)\right)
                         \frac{dk}{(2\pi)^4}\,dv.
 \label{wcr-diagram}
\end{align}
The lattice integral is over the Brillouin zone, the continuum integral
over \(\mathbb R^4\). This identity is obtained from compact-source
Duhamel variations, or defines their absolutely convergent two-point
kernel. In the continuum Gaussian smoothing controls the first-order
differential insertions at each \(s>0\). Only total-trace cyclicity is
used.

For \(r\le2\), direct differentiation or Duhamel gives
\begin{align}
 \|\partial_k^rE(t,k)\|
   &\le C\rho^{-r}(1+t\rho^2)^r e^{-ct\rho^2},\notag\\
 \|\partial_k^r(E_W-E_c)(t,k)\|
   &\le Ct\rho^{4-r}(1+t\rho^2)^r e^{-ct\rho^2}.
 \label{wcr-heat-errors}
\end{align}
These estimates hold at \(t=0\) as well; there is no inverse heat time.

After two external derivatives, the tadpole vertex difference is bounded
by \(C\rho^{-1}e^{-cs\rho^2}\). Its integral with the prefactor \(s\)
is \(O(s^{-1/2})\) in four dimensions.
The heat-difference part is bounded by \(Cs\rho^4e^{-cs\rho^2}\)
before that prefactor, giving \(O(s^{-2})\).

Telescope the bubble product into four differences.
Replacing a vertex gives combined degree \(2+1=3\). Distributing
two external derivatives among all factors therefore gives
\[
 C\rho P(s\rho^2)e^{-cs\rho^2}.
\]
Its integral times \(s^2\) is \(O(s^{-1/2})\).
Replacing a heat factor instead gives
\[
 Cs\rho^4P(s\rho^2)e^{-cs\rho^2},
\]
whose integral times \(s^2\) is \(O(s^{-1})\).
The estimates are uniform in \(v\in[0,1]\): the two heat times sum to
\(s\), and \eqref{wcr-heat-errors} has no endpoint singularity.
Away from the low-momentum ball, both the lattice and continuum
integrals are a polynomial in \(s\) times \(e^{-cs}\).
The continuum tail uses its polynomial vertices and Gaussian decay.
Thus
\begin{equation}
 |\partial_p^2({\cal H}_{W}-{\cal H}_{c})(s,0)|
       \le Cs^{-1/2},\qquad s\ge1.
 \label{wcr-kernel-bound}
\end{equation}
Equivalently one can rescale \(k=\xi/\sqrt{s}\);
the displayed bounds provide the needed common integrable majorant.

The lattice gauge-invariant diagonal heat density and its zero linear
color jet give
\begin{equation}
 h(s)=\tfrac14\partial_{p_\nu}^2
              {\cal H}_{\alpha H,\alpha H}(s,0),\qquad\alpha\ne\nu.
 \label{wcr-density-factor}
\end{equation}
Continuum source kernels may contain contact distributions, so we do
not assume absolute continuum spatial moments or invoke a continuum
reanchoring theorem. The direct transverse diagram calculation below
identifies one quarter of the continuum derivative with its explicit
magnetic density. This and \eqref{wcr-kernel-bound} compare the affine
coefficients with error \(Cs^{-1/2}\).

\subsection{Exact continuum coefficient and sign}
\label{wcr-value-section}

Let \(q\) range over all eigenvalues of \(\operatorname{ad}_H\) on the
complexified adjoint space, with both charge signs and their
multiplicities. For one constant curvature \(F_{12}=fH\), the scalar
diagonal heat kernel in this charge sector is
\begin{equation}
 k_q(s;f)=(4\pi s)^{-2}\frac{qfs}{\sinh(qfs)}.
 \label{wcr-scalar-value}
\end{equation}
Indeed Fourier transformation in \(x_2\) gives oscillators with
eigenvalues \(|qf|(2n+1)\); integration over their centers contributes
\(|qf|/(2\pi)\). Summing the levels and multiplying by the two free
directions proves \eqref{wcr-scalar-value} as a pointwise density.

The spin matrix in \eqref{wcr-continuum-operators} has eigenvalues
\(0,0,2qf,-2qf\), so the vector density is
\[
 k_q(s;f)\{2+2\cosh(2qfs)\}.
\]
Expanding at \(f=0\) for each fixed \(s>0\) gives
\begin{equation}
 [f^2]k_q=-\frac{q^2}{96\pi^2},\qquad
 [f^2](\hbox{vector density})=\frac{10q^2}{48\pi^2}.
 \label{wcr-one-charge}
\end{equation}
This uses \(x/\sinh x=1-x^2/6+O(x^4)\) and
\(2+2\cosh(2x)=4+4x^2+O(x^4)\).
Complexification has the same trace as the original real color space,
so summing charges introduces no extra factor two.
Equations \eqref{wcr-kernel-bound}--\eqref{wcr-one-charge}
prove \eqref{wcr-fixed-bound}.

For the required direct identification of the continuum derivative,
take an external
\(a_\mu\) with momentum \(p=p_\nu e_\nu\), \(\mu\ne\nu\).
The scalar continuum first vertex is \(2qk_\mu\), while its second
vertex \(2q^2\) has no external derivative. The bubble is
\[
 2q^2s(4\pi s)^{-2}\int_0^1e^{-sv(1-v)p^2}\,dv.
\]
Its second \(p\)-derivative is \(-q^2/(24\pi^2)\), giving the scalar
value in \eqref{wcr-one-charge} after the factor \(1/4\).
The vector orbital term is four copies of this scalar term.
The two spin insertions have trace product \(8q^2p^2\), adding
\(q^2/(4\pi^2)\) to the affine coefficient and yielding the vector
value above. Mixed spin/orbital terms have zero Lorentz trace.
No spatial-moment assertion about contact distributions is needed.

For the determinant combination
\(\log\det G-\tfrac12\log\det V\), the proper-time quadratic density
is \(\tfrac12h_v-h_g\). Extraction with \(\beta=2[t^2]\) therefore
gives \(h_v-2h_g\), with this sign and no further multiplier.
Together with the actual-base bound, its middle band is
\(11\alpha_g\log(K^2)+O(1)\), proving \eqref{wcr-middle}.
For \(\tau=\operatorname{Tr}_{\rm fund}/N\) on \(SU(N)\),
\(C_H=2N^2\).

This is a positive-time density comparison. It does not assert the
existence of separate massless infinite-volume determinants at nonzero
magnetic field. Actual soft factors, their joint determinant identities
and bounded outer-band corrections must still be assembled in their
proved domains; no finite-regulator trajectory or continuum existence
is inferred from this reference computation alone.



\section{The actual logarithm and a positive structural margin}
\label{apb-section}

Return to the exactly normalized determinant correction
\(\mathcal D_{M,K}\) of Notebooks 32--36. All stationary substitutions,
coarse constraint factors, and gauge-density factors are those in
that construction. The following conclusion combines estimates for
that object; it is not a replacement of its determinant by a reference.

\begin{theorem}[Uniform actual leading coefficient]
\label{apb-main}
Under the classical structural, determinant-identification and
localized-vertex inputs stated in this notebook and Notebook 36,
there is a constant \(C_{\rm cmp}\), independent of \(M,K\) and of
the applicable joint infinite structural system, such that
\begin{equation}
 \left|\beta(\mathcal D_{M,K})
                    -22\alpha_g\log K\right|\le C_{\rm cmp},
 \qquad \alpha_g=\frac{C_H}{48\pi^2}.
 \label{apb-coefficient}
\end{equation}
The same constant works for every jointly obtainable structural limit
supplying these uniform inputs, including limits with varying old depth.
This is not an exact continuous-momentum identity at finite period.
\end{theorem}

\begin{proof}
Write, with the endpoint-\(K\) affine normalization,
\[
 h^{\rm soft}_{M,K}(s)
  =[t^2]\{\operatorname{Tr}_{0,E}e^{-sV(t)}
                   -2\operatorname{Tr}_{0,S}e^{-sT_s(t)}\}.
\]
The exact soft factorization, its two bounded coarse corrections,
and the two heat-end estimates of Notebook 36 give
\begin{equation}
 \beta(\mathcal D_{M,K})
       =\int_1^{K^2}h^{\rm soft}_{M,K}(s)\,\frac{ds}{s}+R_{M,K},
 \qquad |R_{M,K}|\le C_0.
 \label{apb-band}
\end{equation}
Theorem~\ref{spb-main} removes the soft penalties on this interval:
\[
 |h^{\rm soft}_{M,K}(s)-h^{\rm base}_M(s)|
       \le C_1s/K^2,\qquad 1\le s\le K^2.
\]
The scalar base is independent of \(M\).
Theorem~\ref{abw-main} compares the actual vector base with the
Wilson base with error at most \(C_2/s\). Theorem~\ref{wcr-main}
gives its normalized continuum limit with error at most
\(C_3s^{-1/2}\). Consequently
\begin{equation}
 |h^{\rm soft}_{M,K}(s)-11\alpha_g|
       \le C\{s^{-1/2}+s/K^2\},
       \qquad 1\le s\le K^2.
 \label{apb-plateau}
\end{equation}
Both errors are integrable uniformly against logarithmic time:
\[
 \int_1^{K^2}s^{-1/2}\frac{ds}{s}\le2,\qquad
 \int_1^{K^2}\frac{s}{K^2}\frac{ds}{s}\le1.
\]
The constant part integrates to
\(11\alpha_g\log(K^2)=22\alpha_g\log K\).
Together with \eqref{apb-band}, this proves
\eqref{apb-coefficient}. At \(K=1\) the intermediate integral is
empty and the same bounded-end argument applies.

Every bound used above depends only on the common classical and
localized-vertex constants. No uniqueness or favorable selection of
a structural limit is used. This proves the stated uniform scope.
\end{proof}

Combining Theorem~\ref{sgc-main} with the exact mean-zero identity
\(\mathcal D=\mathcal N_0(\log\det\widehat G
-\tfrac12\log\det Q_B)\) also gives
\[
 \beta(\mathcal N_0\log\det\widehat G)
       =2\alpha_g\log K+O(1),\qquad
 \beta(-\tfrac12\mathcal N_0\log\det Q_B)
       =20\alpha_g\log K+O(1).
\]
These separated values concern this particular exact decomposition.
The complete coefficient \eqref{apb-coefficient} retains the original
normalization.

For a compact simple group and a nonzero normalized Cartan element,
\(C_H>0\). Fixing the elementary quantum ratio \(R>1\), sufficiently
large fixed multiples \(q\) of the activity-return length therefore
give a positive margin for the long ratio \(K=R^q\), uniformly in
all the structural limits required by finite transfer. The next
section states exactly what this supplies for the local quantum
recursion. Neither the present coefficient nor that recursion
identifies the full Wilson measure or constructs continuum fields.


% No preamble. Prefix lpd-. Conditional downstream theorem, 17 October 2026.
% The actual uniform coefficient comparison is a named hypothesis here.
\section{From a uniform positive coefficient to finite local trajectories}
\label{lpd-section}

This section isolates the exact downstream consequence of a uniform
actual determinant comparison. The finite compactness, moment passage,
activity induction and endpoint shooting inputs are those of
Notebooks 31--32. The coefficient comparison below is an explicit
hypothesis of this section; it must be supplied by the actual scalar
and vector comparison, with the stated uniformity. No perturbative
reference coefficient is substituted for an unidentified actual one.

\subsection{Fixed data and the uniform coefficient hypothesis}
\label{lpd-data}

Fix the elementary quantum ratio \(R>1\), grouping \(J\), field
radius \(b>0\), carrier exponent \(a_0\), and the real extracted
finite activity system. Let \(q_a\) be its amplitude block length,
\(r_b=K_{\mathrm{amp}}b^2\) its endpoint ball radius, and
\(C_{\mathrm{int}}b^2\) its intermediate bound. These data are fixed
before the longer block length used below. The activity induction
holds for arbitrary prescribed couplings in one fixed weak range,
provided the retained periods exceed its threshold.

All finite extraction functionals \(\beta_j\) are the actual
short-lift functionals normalized by the actual classical reference.
For the actual finite local map \(F_{j,g}\), write
\begin{equation}
 d_j(g,x)=\beta_{j+1}(F_{j,g}(x)),\qquad
 x_{j+1}=P_{j+1}F_{j,g_j}(x_j),\qquad
 \gamma_{j+1}=\gamma_j-d_j(g_j,x_j).
 \label{lpd-updates}
\end{equation}
Real initial classes and the real branches are understood throughout.
The initial class may be zero, or any fixed real extracted class in
the endpoint ball. The actual coupling relation is
\(g_j=\gamma_j^{-1/2}\) whenever \(\gamma_j>0\).

Let \(D_{M,K}^{*}\) denote the actual leading determinant correction
at preceding homogeneous depth \(M\) and endpoint ratio \(K\) in a
joint infinite structural system. A star labels that structural
system, not a preferred limit of one fixed old depth. Its Maxwell
functional is the true infinite one, with the same classical
normalization used in the finite moment passage. The system includes
the shifted depths and endpoint ratios needed by the determinant
cocycle. In particular all comparisons below refer to the same
joint system.

\paragraph{Uniform actual coefficient comparison \(\mathrm{UC}\).}
Assume that there are \(\alpha_g>0\), \(C_{\mathrm{cmp}}\ge0\)
and \(K_0\ge1\), independent of preceding depth and of the applicable
joint structural limit, such that
\begin{equation}
 \left|\beta^{*}(D_{M,K}^{*})
                  -22\alpha_g\log K\right|
       \le C_{\mathrm{cmp}}
 \quad\hbox{for every }K=R^n\ge K_0.
 \label{lpd-UC}
\end{equation}
For the one-color convention used in the scalar comparison,
\(\alpha_g=C_H/(48\pi^2)\), with \(C_H>0\) for a nonzero probe
in the simple color factor. Only the uniform lower half of
\eqref{lpd-UC} is required for this section. Positivity in one
chosen fixed-depth limit is insufficient.

In the present assembly, Theorem~\ref{apb-main} supplies
\(\mathrm{UC}\), with \(K_0=1\), under its stated classical structural
and determinant inputs. The following argument retains the hypothesis
as a named interface so that its downstream use can be checked
independently.

\subsection{A positive leading block in every structural limit}
\label{lpd-leading-block}

Choose a positive multiple \(q\) of \(q_a\) such that
\begin{equation}
 R^q\ge K_0,\qquad
 C_{\mathrm{cmp}}\le 11\alpha_g q\log R,
 \qquad
 \delta=11\alpha_g q\log R>0.
 \label{lpd-q-choice}
\end{equation}
The exact leading cocycle and true infinite classical marginal
preservation give, in every relevant joint system,
\begin{equation}
 B_q^{*}:=\sum_{j=0}^{q-1}\beta_{j+1}^{*}(D_j^{*})
       =\beta_q^{*}(D_{0,R^q}^{*})
       \ge 22\alpha_g q\log R-C_{\mathrm{cmp}}
       \ge\delta.
 \label{lpd-infinite-block}
\end{equation}
The index zero on the right means the beginning of this block;
its preceding homogeneous depth is arbitrary. The cocycle identity
is used only in the joint infinite system. No exact additivity
of finite short-lift ratios is assumed.

Let \(\mathcal P\) be the compatible cofinal family of periods.
Every fixed power of \(R\), and every required fixed grouping,
divides all sufficiently late periods. For a \(q\)-step block
with final period \(N\), write
\begin{equation}
 M_j=MR^j,\qquad Q_j=R^{q-j}N,\qquad
 b_{0,j}^{M,N}
  =\beta_{M_jR,Q_{j+1}}(D_{M_j,R;Q_{j+1}}),\qquad
 B_q(M,N)=\sum_{j=0}^{q-1}b_{0,j}^{M,N}.
 \label{lpd-finite-block}
\end{equation}
The available joint compactness and normalized moment passage imply
that, along a subsequence of every sequence \(M_n\) and
\(N_n\in\mathcal P\) with \(N_n\to\infty\),
\begin{equation}
 B_q(M_n,N_n)\longrightarrow B_q^{*}.
 \label{lpd-finite-passage}
\end{equation}
This passage includes the long-ratio determinant in
\eqref{lpd-infinite-block}; it does not select unrelated favorable
limits. Exponential source-Hessian bounds control the second-moment
tails, and the actual finite reference denominators tend uniformly
to one.

\begin{lemma}[Uniform positive finite leading block]
\label{lpd-finite-drift}
Under \(\mathrm{UC}\), there is \(N_{\mathrm{drift}}\), independent
of the preceding depth \(M\), such that
\begin{equation}
 B_q(M,N)\ge\delta/2
 \quad\hbox{whenever }N\in\mathcal P,\quad
                         N\ge N_{\mathrm{drift}}.
 \label{lpd-finite-leading-positive}
\end{equation}
\end{lemma}
\begin{proof}
Otherwise choose violating pairs \((M_n,N_n)\) with
\(N_n\to\infty\). Joint compactness and
\eqref{lpd-finite-passage} produce a structural limit with
\(B_q^{*}\le\delta/2\), contradicting
\eqref{lpd-infinite-block}. The starting stage of a block is absorbed
into \(M\). This is a qualitative compactness argument and supplies
no numerical period threshold or convergence rate.
\end{proof}

\subsection{From leading drift to the actual finite decrement}
\label{lpd-actual-drift}

The finite leakage, zero-input parity and inherited-input comparison
give, on the established activity envelope,
\begin{align}
 |d_j(g,x_j)-b_{0,j}|
 \le{}&C_{\mathrm{int}}\omega(N_{\min})
    +C_{\mathrm{par},\beta}
          \{\tau(T(g))+g^2T(g)^6\}\notag\\
 &+C_{\mathrm{in},\beta}\,gT(g)/b .
 \label{lpd-error}
\end{align}
Here \(\omega(N)\to0\) along the compatible family, and the last
two terms tend to zero as \(g\downarrow0\), at the fixed geometric
data and field radius. The cutoff \(T(g)\) is the actual continuous
real threshold. The input factor \(b^2\) has already canceled the
\(b^{-2}\) Maxwell-extraction bound in the last term.

Choose \(N_0\) above \(N_{\mathrm{drift}}\), the reference and
finite-induction thresholds, and sufficiently large that
\begin{equation}
 C_{\mathrm{int}}\omega(N_0)\le\frac{\delta}{8q}.
 \label{lpd-period-threshold}
\end{equation}
Lower the weak-coupling ceiling \(g_{\mathrm w}>0\), if necessary,
so that for every \(0<g\le g_{\mathrm w}\),
\begin{equation}
 C_{\mathrm{par},\beta}
       \{\tau(T(g))+g^2T(g)^6\}
      +C_{\mathrm{in},\beta}gT(g)/b
       \le\frac{\delta}{8q}.
 \label{lpd-coupling-threshold}
\end{equation}
These choices are independent of the permitted history length.
The existing magnitude bound gives a fixed \(D_{\mathrm{abs}}>0\)
with \(|d_j(g,x_j)|\le D_{\mathrm{abs}}\) on this envelope.
Consequently every \(q\)-step block at arbitrary prescribed weak
couplings, started in the endpoint ball and with final period at
least \(N_0\), obeys
\begin{equation}
 B_-\le\sum_{j=k}^{k+q-1}d_j(g_j,x_j)\le B_+,\qquad
 B_-=\delta/4>0,\qquad B_+=qD_{\mathrm{abs}}.
 \label{lpd-block-bounds}
\end{equation}
Each \(q_a\)-step subblock returns to the same endpoint ball.
Thus taking \(q\) larger repeats the existing activity induction;
it does not introduce a new large-endpoint field-dilation estimate.
Neither individual leading coefficients nor individual actual
decrements need be positive.

\subsection{Finite local endpoint shooting}
\label{lpd-shooting}

\begin{theorem}[Conditional positive-drift local trajectories]
\label{lpd-main}
Assume \(\mathrm{UC}\) with the stated structural uniformity and
the established finite local inputs above. Put
\(\Gamma_{\mathrm w}=g_{\mathrm w}^{-2}\), and choose a terminal
inverse coupling
\begin{equation}
                 \Gamma_*\ge
                    \Gamma_{\mathrm w}+qD_{\mathrm{abs}}.
 \label{lpd-terminal}
\end{equation}
For every integer \(m\ge1\) and every compatible finite homogeneous
history of \(mq\) steps with final period \(N_{\mathrm{end}}\ge N_0\),
there is an initial inverse coupling
\begin{equation}
 \Gamma_*+mB_-\le\gamma_{\mathrm{bare}}
                      \le\Gamma_*+mB_+
 \label{lpd-initial-interval}
\end{equation}
whose actual local extraction trajectory reaches \(\Gamma_*\)
and stays in the weak range throughout. Its initial period is
\(R^{mq}N_{\mathrm{end}}\); all intermediate periods obey the
same compatibility requirements. Its activity classes obey the
same endpoint and intermediate \(O(b^2)\) bounds.
\end{theorem}

\begin{proof}
For the proof alone, evaluate the local map at
\begin{equation}
 g_{\mathrm{eval}}(\gamma)
             =\max\{\gamma,\Gamma_{\mathrm w}\}^{-1/2},
 \label{lpd-clipped-coupling}
\end{equation}
while retaining the unclipped scalar update
\(\gamma_{\mathrm{next}}
 =\gamma-d_j(g_{\mathrm{eval}}(\gamma),x)\).
The evaluated couplings are always weak, so finite induction and
\eqref{lpd-block-bounds} apply to every auxiliary history, regardless
of the scalar value. The prescribed-coupling estimates apply to
its realized sequence; no adaptation changes their constants.

For this fixed finite regulator the joint coupling/class continuity
of the normalized local map makes the terminal scalar continuous
in its initial value. On the compact interval
\eqref{lpd-initial-interval}, the absolute increment bound keeps
all intermediate scalars bounded. Their evaluated couplings
therefore lie in a compact positive interval. No continuity at
\(g=0\), volume-uniform continuity modulus, or continuous choice
of equivalent activity representatives is needed.

At the left endpoint of \eqref{lpd-initial-interval}, the terminal
scalar is at most \(\Gamma_*\), by the lower block bound. At its
right endpoint it is at least \(\Gamma_*\), by the upper block
bound. The intermediate value theorem selects a clipped history
ending exactly at \(\Gamma_*\). Summing its future block decrements
backwards gives
\begin{equation}
 \Gamma_*+(m-k)B_-\le\gamma_{kq}
                   \le\Gamma_*+(m-k)B_+,
 \qquad 0\le k\le m.
 \label{lpd-block-endpoints}
\end{equation}
Inside each block,
\begin{equation}
 \gamma_{kq+r}\ge\gamma_{kq}-rD_{\mathrm{abs}}
       \ge\Gamma_*-qD_{\mathrm{abs}}
       \ge\Gamma_{\mathrm w},
 \qquad 0\le r\le q.
 \label{lpd-clipping-absent}
\end{equation}
Clipping is therefore absent from the selected history.
Its actual updates are exactly \eqref{lpd-updates} with
\(g_j=\gamma_j^{-1/2}\). The finite quotient relocalizations
replace representatives only by equal normalized totals; their
scalar differences are carried by the exact vacuum factors.
Thus the conclusion concerns the genuine local density update,
not an approximate scalar recurrence.
\end{proof}

The bare coupling in this theorem satisfies
\begin{equation}
 (\Gamma_*+mqD_{\mathrm{abs}})^{-1/2}
       \le g_{\mathrm{bare}}
       \le(\Gamma_*+m\delta/4)^{-1/2}.
 \label{lpd-bare-asymptotic}
\end{equation}
It is of order \(m^{-1/2}\), equivalently of order
\((mq)^{-1/2}\) with \(q\) fixed, and reaches the same positive
weak endpoint \(g_*=\Gamma_*^{-1/2}\). The theorem asserts neither
uniqueness of the tuned initial value nor consistency of the
choices at different horizons.

\subsection{Parameter order and scope}
\label{lpd-scope}

First fix \(R,J,b\) and the available \(q_a\)-block activity system.
The uniform coefficient comparison determines a sufficiently large
fixed multiple \(q\) by \eqref{lpd-q-choice}. Then choose the
period threshold \(N_0\), the weak ceiling \(g_{\mathrm w}\), and
the fixed terminal margin \(\Gamma_*\). The horizon \(m\) is
arbitrary after these choices. No quantitative finite-volume rate
and no exact finite classical marginal identity enter this order.

The only new coefficient input of this section is
\eqref{lpd-UC}. If an actual comparison theorem supplies it in
every required joint structural system, this section discharges
the positive-drift hypothesis of finite local shooting. The
comparison is supplied here by Theorem~\ref{apb-main}. The
conclusion remains about homogeneous local analytic histories
and their justified continuation. It does not identify those
histories with the full Wilson measure, control mixed or large-field
sectors, construct continuum observable measures, prove a
nontrivial continuum Yang--Mills theory, or establish a mass gap.


\section{What must still become a four-dimensional quantum theory}
\label{pdh-existence-boundary}

The new coefficient removes a local ultraviolet obstruction. It does
not turn the analytic local recursion into the full Wilson measure.
The next construction must identify that recursion, with its retained
windows, gauge insertion, coupling extraction and vacuum convention,
inside an exact decomposition of the original finite Wilson integral.
Every mixed and large-field history must be included and estimated
with compatible boundary conditions. An arbitrary continuation of a
small-field germ outside its domain would define a different model.

The next priority is therefore a full-density matching theorem for
one fixed scheme. A useful statement would start from the actual
positive finite Wilson measure and express its pushforward as the
proved local sector plus explicitly controlled complementary sectors,
at the tuned couplings just constructed. It must hold uniformly as
the ultraviolet horizon grows. This statement is a target, not a
consequence of the one-loop calculation.

External observable sources then require their own uniform estimates.
Holomorphy in a connection background does not establish holomorphy
in a source coupled to a renormalized observable. Source counterterms,
normalization denominators, reflection-cell domains, and estimates
of all needed orders must be supplied for the same full measure.
Those estimates must produce a common continuum hierarchy with
compatible products, rather than separate limits of unrelated
observables.

Reflection positivity must pass from that same positive Wilson
sequence. Full Euclidean covariance, the required distributional
growth, a suitable vacuum sector, and reconstruction must be verified.
Finally, a nonzero scalar correlation alone does not identify
Yang--Mills: the limit needs an adequate mixed family of gauge-invariant
curvature observables and their short-distance normalization and
transformation laws. An interacting lower bound and the relevant
identification remain separate tasks. The mass-gap problem requires
an additional physical infrared argument.

The companion file
\texttt{work/final\_push\_existence\_dependency\_audit.txt}
records the detailed dependency map and source boundaries.
It explicitly distinguishes established working modules, conditional
reductions, and the open model-specific inputs. Earlier notebooks
remain chronological records; their older next-step paragraphs
must be read together with the results of the present notebook.

\begin{thebibliography}{9}
\bibitem{prior37}
\emph{Toward a four-dimensional Yang--Mills construction},
research Notebooks 15--36 in this workspace. In particular,
Notebook 32 supplies finite moment transfer and local shooting;
Notebook 33 fixes the continuum normalization and affine density
identification; Notebook 34 computes the exact unpinned scalar heat
coefficient; Notebooks 35--36 supply the actual mean-zero and soft
determinant identities and localized estimates.
\bibitem{vassilevich37}
D.~V.~Vassilevich, \emph{Heat kernel expansion: user's manual},
Physics Reports \textbf{388} (2003), 279--360,
\href{https://arxiv.org/abs/hep-th/0306138}{arXiv:hep-th/0306138}.
The reference coefficient is calculated explicitly here;
the source is not a proof of the actual lattice comparison.
\bibitem{adamslee37}
D.~H.~Adams and W.~Lee,
\emph{Structure of logarithmically divergent one-loop lattice
Feynman integrals}, Phys.\ Rev.\ D \textbf{77} (2008), 045010,
\href{https://arxiv.org/abs/0709.0781}{arXiv:0709.0781}.
Its fixed-integrand universality theorem is useful background;
uniformity in our varying effective actions is proved directly above.
\bibitem{jw37}
A.~Jaffe and E.~Witten, \emph{Quantum Yang--Mills Theory},
\href{https://www.claymath.org/wp-content/uploads/2022/06/yangmills.pdf}
{Clay Mathematics Institute problem description}.
\bibitem{source37}
E.~Sussman and collaborators,
\href{https://github.com/ethanwsussman/Yang--Mills-LLM-experiment}
{Yang--Mills-LLM-experiment}, deposited accounts pinned at commit
\texttt{e392ef4501\allowbreak 3a581664ba\allowbreak b76b908671\allowbreak 70b30fe10e}.
\end{thebibliography}
\end{document}
